\documentclass[paper=a4,fontsize=11pt]{jlreq}

% 必要なパッケージ
\usepackage{luatexja} % LuaLaTeXで日本語を使うため
\usepackage{listings} % ソースコード表示用
\usepackage{xcolor}   % ソースコードの色付け用

% ソースコード表示(listings)の設定
\lstset{
    basicstyle=\ttfamily\small, % 基本フォント
    numbers=left,               % 行番号を左に表示 (課題要件)
    numberstyle=\tiny\color{gray}, % 行番号のスタイル
    frame=single,               % コードを枠線で囲む
    tabsize=4,                  % タブ幅
    breaklines=true,            % 自動改行
    captionpos=t,               % キャプションの位置(上)
    showstringspaces=false      % 文字列中の空白を記号にしない
}

\title{ソフトウェア演習2B 第1回レポート}
\author{学籍番号: 00000000 \quad 氏名: 豊橋 太郎}
\date{\today}

\begin{document}

\maketitle

% --- 課題1 ---
\section{課題1}
\subsection{課題内容}
% ここに課題1の内容を記述します。
（課題内容を記述）

\subsection{ソースコード}
% 言語を指定してソースコードを貼り付けます（例: C++）
\begin{lstlisting}[language=C++, caption=Message.cpp]
#include <iostream>
#include "Message.h"

int main (int argc, char *argv[]) {
    Message obj;
    obj.setMessage("Hello World.");
    std::cout << obj.getMessage() << std::endl;
    return 0;
}
\end{lstlisting}

\subsection{プログラムに関する説明}
% プログラムの解説やアピールポイントを記述します。
（作成したプログラムに関する説明、工夫した点、アピールポイントなどを記述）


% --- 課題2 ---
\section{課題2}
\subsection{課題内容}
（課題内容を記述）

\subsection{ソースコード}
\begin{lstlisting}[language=C++, caption=main.cpp (課題2)]
// 課題2のソースコードを記述
\end{lstlisting}

\subsection{プログラムに関する説明}
（作成したプログラムに関する説明を記述）


% --- 課題3 ---
\section{課題3}
\subsection{課題内容}
（課題内容を記述）

\subsection{ソースコード}
\begin{lstlisting}[language=C++, caption=main.cpp (課題3)]
// 課題3のソースコードを記述
\end{lstlisting}

\subsection{プログラムに関する説明}
（作成したプログラムに関する説明を記述）


% --- 課題4 (チェック項目) ---
\section{課題4 チェック項目への回答}
% 課題資料の最後にあるチェック項目への回答を記述します。
\begin{enumerate}
    \item \textbf{private, protected, publicの違いについて} \\
    （ここへ回答を記述します）
    \item \textbf{その他のチェック項目} \\
    （ここへ回答を記述します）
\end{enumerate}

\end{document}